\chapter{拉格朗日方程与能量法}

\intro{
	最完美的方程是那些不仅具有真理性，而且具有美感的方程。
	
	\rightline{——Paul Dirac}
}

\section{分析力学基础}

Newton 力学以矢量形式处理力与加速度，而\textbf{分析力学}以能量（标量）为基础，提供了一种更系统化的建模框架，尤其适用于复杂约束系统。其核心优势在于：
\begin{enumerate}
	\item 自动消除约束反力，仅需处理主动力；
	\item 以标量函数（动能 $T$、势能 $V$）为出发点，推导过程系统化；
	\item 形式与坐标系选择无关，便于数值实现。
\end{enumerate}

\subsection{约束的分类与自由度}

力学系统的运动通常受到约束的限制。约束可以从多个角度分类：

\begin{mydefinition}{完整约束与非完整约束}
	\textbf{完整约束}（Holonomic constraint）可表示为位形和时间的等式：
	\[
	f(\mathbf{r}_1, \mathbf{r}_2, \ldots, \mathbf{r}_N, t) = 0
	\]
	或等价地用广义坐标表示：$f(q_1, \ldots, q_n, t) = 0$。\\
	\textbf{非完整约束}（Nonholonomic constraint）不可积分为位形空间的等式，通常涉及速度：
	\[
	\sum_{j} a_{j}(q, t) \dot{q}_j + b(q, t) = 0
	\]
\end{mydefinition}

\begin{mydefinition}{定常约束与非定常约束}
	\textbf{定常约束}（Scleronomic）：约束方程 $f(\mathbf{r}_i) = 0$ 不显含时间 $t$。\\
	\textbf{非定常约束}（Rheonomic）：约束方程 $f(\mathbf{r}_i, t) = 0$ 显含时间 $t$。
\end{mydefinition}

\begin{example}{约束分类实例}
	单摆——绳长 $l$ 不变：$x^2 + y^2 - l^2 = 0$（完整、定常约束）。\\
	冰刀在冰面上滑动——速度方向沿刀口方向：$\dot{x} \sin\theta - \dot{y} \cos\theta = 0$（非完整约束）。
\end{example}

\subsection{广义坐标与自由度}

\begin{mydefinition}{广义坐标与自由度}
	能够完全确定系统位形的一组独立参数 $q_1, q_2, \ldots, q_n$ 称为\textbf{广义坐标}。\\
	对于含 $m$ 个完整约束的 $N$ 个质点系统：\textbf{自由度} $n = 3N - m$。
\end{mydefinition}

位矢可表示为广义坐标（与时间）的函数：
\[
\mathbf{r}_i = \mathbf{r}_i(q_1, q_2, \ldots, q_n, t), \quad i = 1, 2, \ldots, N
\]
速度由链式法则给出：
\[
\dot{\mathbf{r}}_i = \sum_{j=1}^{n} \frac{\partial \mathbf{r}_i}{\partial q_j} \dot{q}_j + \frac{\partial \mathbf{r}_i}{\partial t}
\]

\subsection{虚位移与虚功}

\begin{mydefinition}{虚位移}
	在给定瞬时 $t$，在约束允许范围内假想的无限小位移 $\delta \mathbf{r}_i$ 称为\textbf{虚位移}。虚位移与实位移的区别在于：
	\begin{itemize}
		\item 虚位移 $\delta \mathbf{r}_i$：瞬时假想位移，与时间进程无关（$\delta t = 0$）；
		\item 实位移 $d\mathbf{r}_i$：在时间增量 $dt$ 内实际发生的位移。
	\end{itemize}
	对于定常约束，虚位移与实位移方向相同；对于非定常约束，两者方向不同。
\end{mydefinition}

虚位移用广义坐标的变分表示为：
\[
\delta \mathbf{r}_i = \sum_{j=1}^{n} \frac{\partial \mathbf{r}_i}{\partial q_j} \delta q_j
\]
注意此处不包含 $\frac{\partial \mathbf{r}_i}{\partial t}$ 项，因为虚位移发生在固定时刻 $(\delta t = 0)$。

\subsection{虚功原理与 D'Alembert 原理}

\begin{theorem}{虚功原理（静力学）}
	受理想约束的平衡力学系统，所有主动力在任意虚位移上所作的虚功之和为零：
	\[
	\delta W = \sum_{i} \mathbf{F}_i^{\text{(a)}} \cdot \delta \mathbf{r}_i = 0
	\]
	\textbf{理想约束}的定义：约束反力 $\mathbf{R}_i$ 在任意虚位移上不作功，即 $\sum_i \mathbf{R}_i \cdot \delta \mathbf{r}_i = 0$（如光滑铰链、刚性杆、不可伸长绳索等）。
\end{theorem}

\begin{proof}[虚功原理的推导]
	设系统处于平衡状态。对每个质点 $i$，合力为零：$\mathbf{F}_i^{\text{(a)}} + \mathbf{R}_i = 0$。
	将每个平衡方程点乘对应的虚位移 $\delta \mathbf{r}_i$ 并求和：
	\[
	\sum_i \left( \mathbf{F}_i^{\text{(a)}} + \mathbf{R}_i \right) \cdot \delta \mathbf{r}_i = 0
	\]
	由理想约束条件 $\sum_i \mathbf{R}_i \cdot \delta \mathbf{r}_i = 0$，立即得：
	\[
	\sum_i \mathbf{F}_i^{\text{(a)}} \cdot \delta \mathbf{r}_i = 0
	\]
	这本质上是将矢量形式的平衡方程投影到约束允许的运动方向上，从而消去了未知的约束反力。
\end{proof}

\begin{theorem}{D'Alembert 原理（动力学）}
	将 Newton 第二定律改写为 $\mathbf{F}_i - m_i \ddot{\mathbf{r}}_i = 0$，引入\textbf{惯性力}（又称 D'Alembert 惯性力）$-m_i \ddot{\mathbf{r}}_i$，动力学问题在形式上转化为静力学平衡问题：
	\[
	\sum_{i} (\mathbf{F}_i - m_i \ddot{\mathbf{r}}_i) \cdot \delta \mathbf{r}_i = 0
	\]
\end{theorem}

\begin{proof}[D'Alembert 原理的推导]
	由 Newton 第二定律对每个质点 $i$：
	\[
	m_i \ddot{\mathbf{r}}_i = \mathbf{F}_i^{\text{(a)}} + \mathbf{R}_i
	\]
	移项得：
	\[
	\mathbf{F}_i^{\text{(a)}} + \mathbf{R}_i - m_i \ddot{\mathbf{r}}_i = 0
	\]
	形式上与静力学平衡方程 $\mathbf{F}_i^{\text{(a)}} + \mathbf{R}_i = 0$ 相同，仅多了一项 $-m_i \ddot{\mathbf{r}}_i$（称为惯性力）。
	
	对每个方程点乘虚位移 $\delta \mathbf{r}_i$ 并求和：
	\[
	\sum_i \left( \mathbf{F}_i^{\text{(a)}} + \mathbf{R}_i - m_i \ddot{\mathbf{r}}_i \right) \cdot \delta \mathbf{r}_i = 0
	\]
	由理想约束条件 $\sum_i \mathbf{R}_i \cdot \delta \mathbf{r}_i = 0$，得：
	\[
	\sum_i \left( \mathbf{F}_i^{\text{(a)}} - m_i \ddot{\mathbf{r}}_i \right) \cdot \delta \mathbf{r}_i = 0
	\]
	为简洁起见，以下将 $\mathbf{F}_i^{\text{(a)}}$ 简记为 $\mathbf{F}_i$。
\end{proof}

\begin{remark}
	D'Alembert 原理的精妙之处在于：它将动力学问题与静力学的虚功原理统一起来，避免了直接处理约束反力，为 Lagrange 方程的导出铺平了道路。物理本质上，它是 Newton 第二定律向约束系统的一个弱形式推广。
\end{remark}

\section{第二类 Lagrange 方程的完整推导}

\subsection{准备工作：两个关键恒等式}

在正式推导之前，先建立两个极为重要的运动学恒等式。由链式法则：
\[
\dot{\mathbf{r}}_i = \sum_{k=1}^{n} \frac{\partial \mathbf{r}_i}{\partial q_k} \dot{q}_k + \frac{\partial \mathbf{r}_i}{\partial t}
\]

\textbf{恒等式 1}（速度对广义速度的偏导数）：
将上式两端对 $\dot{q}_j$ 求偏导，由于 $\frac{\partial \mathbf{r}_i}{\partial q_k}$ 与 $\frac{\partial \mathbf{r}_i}{\partial t}$ 均不显含 $\dot{q}_j$：
\[
\frac{\partial \dot{\mathbf{r}}_i}{\partial \dot{q}_j} = \frac{\partial \mathbf{r}_i}{\partial q_j}
\tag{1}
\]

\textbf{恒等式 2}（速度对广义坐标的全导数与偏导数的交换）：
\[
\frac{d}{dt} \left( \frac{\partial \mathbf{r}_i}{\partial q_j} \right) = \sum_{k=1}^{n} \frac{\partial^2 \mathbf{r}_i}{\partial q_k \partial q_j} \dot{q}_k + \frac{\partial^2 \mathbf{r}_i}{\partial t \partial q_j}
\]
另一方面，$\dot{\mathbf{r}}_i$ 对 $q_j$ 求偏导：
\[
\frac{\partial \dot{\mathbf{r}}_i}{\partial q_j} = \frac{\partial}{\partial q_j} \left( \sum_{k=1}^{n} \frac{\partial \mathbf{r}_i}{\partial q_k} \dot{q}_k + \frac{\partial \mathbf{r}_i}{\partial t} \right) = \sum_{k=1}^{n} \frac{\partial^2 \mathbf{r}_i}{\partial q_j \partial q_k} \dot{q}_k + \frac{\partial^2 \mathbf{r}_i}{\partial q_j \partial t}
\]
假设 $\mathbf{r}_i$ 二阶连续可微，混合偏导可交换次序，则两者相等：
\[
\frac{d}{dt} \left( \frac{\partial \mathbf{r}_i}{\partial q_j} \right) = \frac{\partial \dot{\mathbf{r}}_i}{\partial q_j}
\tag{2}
\]

\subsection{从 D'Alembert 原理到 Lagrange 方程}

\begin{proof}[Lagrange 方程的完整推导]
	
	\textbf{Step 1 —— 代入虚位移表达式。}
	
	D'Alembert 原理：
	\[
	\sum_{i} \left( \mathbf{F}_i - m_i \ddot{\mathbf{r}}_i \right) \cdot \delta \mathbf{r}_i = 0
	\]
	代入 $\delta \mathbf{r}_i = \sum_{j} \frac{\partial \mathbf{r}_i}{\partial q_j} \delta q_j$：
	\[
	\sum_{i} \left( \mathbf{F}_i - m_i \ddot{\mathbf{r}}_i \right) \cdot \sum_{j=1}^{n} \frac{\partial \mathbf{r}_i}{\partial q_j} \delta q_j = 0
	\]
	
	\textbf{Step 2 —— 交换求和次序，定义广义力。}
	
	交换对 $i$ 和 $j$ 的求和次序：
	\[
	\sum_{j=1}^{n} \left[ \sum_{i} \mathbf{F}_i \cdot \frac{\partial \mathbf{r}_i}{\partial q_j} - \sum_{i} m_i \ddot{\mathbf{r}}_i \cdot \frac{\partial \mathbf{r}_i}{\partial q_j} \right] \delta q_j = 0
	\]
	
	定义\textbf{广义力}：
	\[
	Q_j \equiv \sum_{i} \mathbf{F}_i \cdot \frac{\partial \mathbf{r}_i}{\partial q_j}
	\tag{3}
	\]
	$Q_j \delta q_j$ 在量纲上为功，广义力的量纲与广义坐标配对：若 $q_j$ 为长度则 $Q_j$ 为力，若 $q_j$ 为角度则 $Q_j$ 为力矩。
	
	由于各 $\delta q_j$ 相互独立（广义坐标的独立性），其系数必须分别等于零：
	\[
	\sum_{i} m_i \ddot{\mathbf{r}}_i \cdot \frac{\partial \mathbf{r}_i}{\partial q_j} = Q_j, \quad j = 1, 2, \ldots, n
	\tag{4}
	\]
	
	\textbf{Step 3 —— 惯性项的恒等变形（推导的核心步骤）。}
	
	现在处理左端的惯性项。考虑单个质点 $i$ 的贡献：
	\[
	m_i \ddot{\mathbf{r}}_i \cdot \frac{\partial \mathbf{r}_i}{\partial q_j}
	\]
	
	利用乘积求导法则：
	\[
	\frac{d}{dt} \left( m_i \dot{\mathbf{r}}_i \cdot \frac{\partial \mathbf{r}_i}{\partial q_j} \right) = m_i \ddot{\mathbf{r}}_i \cdot \frac{\partial \mathbf{r}_i}{\partial q_j} + m_i \dot{\mathbf{r}}_i \cdot \frac{d}{dt} \left( \frac{\partial \mathbf{r}_i}{\partial q_j} \right)
	\]
	移项得：
	\[
	m_i \ddot{\mathbf{r}}_i \cdot \frac{\partial \mathbf{r}_i}{\partial q_j} = \frac{d}{dt} \left( m_i \dot{\mathbf{r}}_i \cdot \frac{\partial \mathbf{r}_i}{\partial q_j} \right) - m_i \dot{\mathbf{r}}_i \cdot \frac{d}{dt} \left( \frac{\partial \mathbf{r}_i}{\partial q_j} \right)
	\tag{5}
	\]
	
	应用恒等式 (1) 和 (2)：
	\[
	\frac{\partial \mathbf{r}_i}{\partial q_j} = \frac{\partial \dot{\mathbf{r}}_i}{\partial \dot{q}_j}, \qquad
	\frac{d}{dt} \left( \frac{\partial \mathbf{r}_i}{\partial q_j} \right) = \frac{\partial \dot{\mathbf{r}}_i}{\partial q_j}
	\]
	
	将恒等式代入 (5)：
	\[
	\begin{aligned}
	m_i \ddot{\mathbf{r}}_i \cdot \frac{\partial \mathbf{r}_i}{\partial q_j} &= \frac{d}{dt} \left( m_i \dot{\mathbf{r}}_i \cdot \frac{\partial \dot{\mathbf{r}}_i}{\partial \dot{q}_j} \right) - m_i \dot{\mathbf{r}}_i \cdot \frac{\partial \dot{\mathbf{r}}_i}{\partial q_j} \quad \leftarrow \text{代入恒等式 (1)} \\
	&= \frac{d}{dt} \left( m_i \dot{\mathbf{r}}_i \cdot \frac{\partial \dot{\mathbf{r}}_i}{\partial \dot{q}_j} \right) - m_i \dot{\mathbf{r}}_i \cdot \frac{\partial \dot{\mathbf{r}}_i}{\partial q_j} \quad \leftarrow \text{代入恒等式 (2)}
	\end{aligned}
	\]
	
	\textbf{关键观察}：注意 $m_i \dot{\mathbf{r}}_i \cdot \frac{\partial \dot{\mathbf{r}}_i}{\partial \dot{q}_j} = \frac{\partial}{\partial \dot{q}_j} \left( \frac{1}{2} m_i |\dot{\mathbf{r}}_i|^2 \right)$，同理 $m_i \dot{\mathbf{r}}_i \cdot \frac{\partial \dot{\mathbf{r}}_i}{\partial q_j} = \frac{\partial}{\partial q_j} \left( \frac{1}{2} m_i |\dot{\mathbf{r}}_i|^2 \right)$。
	
	这是因为：
	\[
	\frac{\partial}{\partial \dot{q}_j} \left( \frac{1}{2} m_i \dot{\mathbf{r}}_i \cdot \dot{\mathbf{r}}_i \right) = \frac{1}{2} m_i \left( \frac{\partial \dot{\mathbf{r}}_i}{\partial \dot{q}_j} \cdot \dot{\mathbf{r}}_i + \dot{\mathbf{r}}_i \cdot \frac{\partial \dot{\mathbf{r}}_i}{\partial \dot{q}_j} \right) = m_i \dot{\mathbf{r}}_i \cdot \frac{\partial \dot{\mathbf{r}}_i}{\partial \dot{q}_j}
	\]
	
	定义第 $i$ 个质点的动能 $T_i = \frac{1}{2} m_i |\dot{\mathbf{r}}_i|^2$，则：
	\[
	m_i \ddot{\mathbf{r}}_i \cdot \frac{\partial \mathbf{r}_i}{\partial q_j} = \frac{d}{dt} \left( \frac{\partial T_i}{\partial \dot{q}_j} \right) - \frac{\partial T_i}{\partial q_j}
	\tag{6}
	\]
	
	\textbf{Step 4 —— 对全部质点求和，得到 Lagrange 方程。}
	
	将 (6) 代入 (4) 并对所有质点求和：
	\[
	\sum_i \left[ \frac{d}{dt} \left( \frac{\partial T_i}{\partial \dot{q}_j} \right) - \frac{\partial T_i}{\partial q_j} \right] = Q_j
	\]
	
	定义系统总动能 $T = \sum_i T_i$，求导与求和可交换次序：
	\[
	\frac{d}{dt} \left( \frac{\partial T}{\partial \dot{q}_j} \right) - \frac{\partial T}{\partial q_j} = Q_j, \quad j = 1, 2, \ldots, n
	\tag{7}
	\]
	
	这就是\textbf{以动能表示的第二类 Lagrange 方程}。
	
	\textbf{Step 5 —— 引入势能，得标准形式。}
	
	将广义力分解为有势力（保守力）与非保守力两部分：
	\[
	Q_j = Q_j^{\text{(cons)}} + Q_j^{\text{(nc)}}
	\]
	
	对于有势力，存在势能函数 $V = V(q_1, \ldots, q_n, t)$ 使得：
	\[
	Q_j^{\text{(cons)}} = -\frac{\partial V}{\partial q_j}
	\]
	且 $V$ 不依赖于广义速度，故 $\partial V / \partial \dot{q}_j = 0$。
	
	定义\textbf{Lagrange 函数}（亦称 Lagrangian）：
	\[
	L(q, \dot{q}, t) = T(q, \dot{q}, t) - V(q, t)
	\]
	
	将 $Q_j$ 代入 (7) 式：
	\[
	\frac{d}{dt} \left( \frac{\partial T}{\partial \dot{q}_j} \right) - \frac{\partial T}{\partial q_j} = -\frac{\partial V}{\partial q_j} + Q_j^{\text{(nc)}}
	\]
	即：
	\[
	\frac{d}{dt} \left( \frac{\partial (T - V)}{\partial \dot{q}_j} \right) - \frac{\partial (T - V)}{\partial q_j} = Q_j^{\text{(nc)}}
	\]
	
	最终得到\textbf{第二类 Lagrange 方程的标准形式}：
	\[
	\boxed{\frac{d}{dt} \left( \frac{\partial L}{\partial \dot{q}_j} \right) - \frac{\partial L}{\partial q_j} = Q_j^{\text{(nc)}}, \quad j = 1, 2, \ldots, n}
	\tag{8}
	\]
\end{proof}

\subsection{Lagrange 方程的结构与物理意义}

Lagrange 方程可写成矩阵形式，揭示其与 Newton 力学中质量-阻尼-刚度体系的结构对应：

\begin{corollary}{线性系统的 Lagrange 方程结构}
	对于在平衡位置附近作微振动的系统，动能和势能可写成二次型：
	\[
	T = \frac{1}{2} \dot{\mathbf{q}}^T \mathbf{M} \dot{\mathbf{q}}, \qquad V = \frac{1}{2} \mathbf{q}^T \mathbf{K} \mathbf{q}
	\]
	其中 $\mathbf{M}$ 为\textbf{质量矩阵}（正定对称），$\mathbf{K}$ 为\textbf{刚度矩阵}（半正定对称）。
	
	代入 Lagrange 方程 (8) 得线性运动方程：
	\[
	\mathbf{M} \ddot{\mathbf{q}} + \mathbf{C} \dot{\mathbf{q}} + \mathbf{K} \mathbf{q} = \mathbf{f}(t)
	\]
	其中 $\mathbf{C}$ 来自 Rayleigh 耗散函数，$\mathbf{f}(t)$ 为外激励广义力。这正是前续章节中反复出现的二阶振动方程的基本形式。
\end{corollary}

\section{Lagrange 方程的应用实例}

在实际建模中，可按以下步骤系统化地应用 Lagrange 方法：
\begin{enumerate}
	\item 选择合适的广义坐标 $q_1, \ldots, q_n$；
	\item 写出各质点位矢与广义坐标的关系，计算速度 $\dot{\mathbf{r}}_i$；
	\item 计算系统动能 $T(q, \dot{q}, t)$ 与势能 $V(q, t)$；
	\item 构造 Lagrange 函数 $L = T - V$；
	\item 计算各偏导数 $\frac{\partial L}{\partial q_j}$、$\frac{\partial L}{\partial \dot{q}_j}$；
	\item 代入 Lagrange 方程 (8)，得到运动微分方程。
\end{enumerate}

\begin{example}{简单摆}
	长度为 $l$ 的单摆。广义坐标取摆角 $\theta$（偏离竖直方向）。
	
	\textbf{步骤 1——运动学}：质点坐标
	\[
	x = l \sin\theta, \quad y = -l \cos\theta
	\]
	速度分量：
	\[
	\dot{x} = l \dot{\theta} \cos\theta, \quad \dot{y} = l \dot{\theta} \sin\theta
	\]
	速率平方：$|\dot{\mathbf{r}}|^2 = \dot{x}^2 + \dot{y}^2 = l^2 \dot{\theta}^2$。
	
	\textbf{步骤 2——能量}：
	\begin{align*}
		T &= \frac{1}{2} m l^2 \dot{\theta}^2 \\
		V &= mgy = -mgl \cos\theta
	\end{align*}
	
	\textbf{步骤 3——Lagrange 函数}：
	\[
	L = T - V = \frac{1}{2} m l^2 \dot{\theta}^2 + mgl \cos\theta
	\]
	
	\textbf{步骤 4——代入 Lagrange 方程}：
	\begin{align*}
		\frac{\partial L}{\partial \dot{\theta}} &= m l^2 \dot{\theta} \\
		\frac{d}{dt} \left( \frac{\partial L}{\partial \dot{\theta}} \right) &= m l^2 \ddot{\theta} \\
		\frac{\partial L}{\partial \theta} &= -mgl \sin\theta
	\end{align*}
	
	Lagrange 方程 $\frac{d}{dt} \left( \frac{\partial L}{\partial \dot{\theta}} \right) - \frac{\partial L}{\partial \theta} = 0$ 给出：
	\[
	m l^2 \ddot{\theta} + mgl \sin\theta = 0 \;\Longrightarrow\; \ddot{\theta} + \frac{g}{l} \sin\theta = 0
	\]
	
	\textbf{线性化}：小角度近似 $\sin\theta \approx \theta$ 下，
	\[
	\ddot{\theta} + \omega_n^2 \theta = 0, \quad \omega_n = \sqrt{\frac{g}{l}}
	\]
	固有周期 $T_n = 2\pi \sqrt{l/g}$，与振幅无关（小角度）。
\end{example}

\begin{example}{弹簧-质量系统}
	质量为 $m$ 的物体连接于刚度系数为 $k$ 的弹簧，水平光滑面上运动。广义坐标取位移 $x$。
	
	\[
	T = \frac{1}{2} m \dot{x}^2, \quad V = \frac{1}{2} k x^2, \quad L = \frac{1}{2} m \dot{x}^2 - \frac{1}{2} k x^2
	\]
	
	Lagrange 方程：
	\[
	\frac{d}{dt} \left( m \dot{x} \right) - (-k x) = 0 \;\Longrightarrow\; m \ddot{x} + k x = 0
	\]
	直接得到标准简谐振动方程 $\ddot{x} + \omega_n^2 x = 0$，$\omega_n = \sqrt{k/m}$。
\end{example}

\begin{example}{双摆系统的详细推导}
	两杆长分别为 $l_1, l_2$，端部质量 $m_1, m_2$（轻杆）。取 $\theta_1, \theta_2$ 为广义坐标。
	
	\textbf{质点 1 的运动学}：
	\begin{align*}
		x_1 &= l_1 \sin\theta_1, \quad y_1 = -l_1 \cos\theta_1 \\
		\dot{x}_1 &= l_1 \dot{\theta}_1 \cos\theta_1, \quad \dot{y}_1 = l_1 \dot{\theta}_1 \sin\theta_1 \\
		|\dot{\mathbf{r}}_1|^2 &= l_1^2 \dot{\theta}_1^2
	\end{align*}
	
	\textbf{质点 2 的运动学}：
	\begin{align*}
		x_2 &= l_1 \sin\theta_1 + l_2 \sin\theta_2 \\
		y_2 &= -l_1 \cos\theta_1 - l_2 \cos\theta_2 \\[3pt]
		\dot{x}_2 &= l_1 \dot{\theta}_1 \cos\theta_1 + l_2 \dot{\theta}_2 \cos\theta_2 \\
		\dot{y}_2 &= l_1 \dot{\theta}_1 \sin\theta_1 + l_2 \dot{\theta}_2 \sin\theta_2
	\end{align*}
	
	速率平方：
	\[
	\begin{aligned}
	|\dot{\mathbf{r}}_2|^2 &= (l_1 \dot{\theta}_1 \cos\theta_1 + l_2 \dot{\theta}_2 \cos\theta_2)^2 + (l_1 \dot{\theta}_1 \sin\theta_1 + l_2 \dot{\theta}_2 \sin\theta_2)^2 \\
	&= l_1^2 \dot{\theta}_1^2 + l_2^2 \dot{\theta}_2^2 + 2 l_1 l_2 \dot{\theta}_1 \dot{\theta}_2 (\cos\theta_1 \cos\theta_2 + \sin\theta_1 \sin\theta_2) \\
	&= l_1^2 \dot{\theta}_1^2 + l_2^2 \dot{\theta}_2^2 + 2 l_1 l_2 \dot{\theta}_1 \dot{\theta}_2 \cos(\theta_1 - \theta_2)
	\end{aligned}
	\]
	
	\textbf{动能}：
	\[
	T = \frac{1}{2} m_1 |\dot{\mathbf{r}}_1|^2 + \frac{1}{2} m_2 |\dot{\mathbf{r}}_2|^2
	\]
	\[
	\boxed{
	T = \frac{1}{2} (m_1 + m_2) l_1^2 \dot{\theta}_1^2 + \frac{1}{2} m_2 l_2^2 \dot{\theta}_2^2 + m_2 l_1 l_2 \dot{\theta}_1 \dot{\theta}_2 \cos(\theta_1 - \theta_2)
	}
	\]
	
	\textbf{势能}（取悬点为势能零点）：
	\[
	V = m_1 g y_1 + m_2 g y_2 = -(m_1 + m_2) g l_1 \cos\theta_1 - m_2 g l_2 \cos\theta_2
	\]
	
	\textbf{Lagrange 函数} $L = T - V$：
	\[
	L = \frac{1}{2}(m_1+m_2)l_1^2 \dot{\theta}_1^2 + \frac{1}{2}m_2 l_2^2 \dot{\theta}_2^2 + m_2 l_1 l_2 \dot{\theta}_1 \dot{\theta}_2 \cos(\theta_1-\theta_2) + (m_1+m_2)gl_1\cos\theta_1 + m_2 g l_2 \cos\theta_2
	\]
	
	\textbf{Lagrange 方程（对 $\theta_1$）}：
	\[
	\frac{\partial L}{\partial \dot{\theta}_1} = (m_1+m_2)l_1^2 \dot{\theta}_1 + m_2 l_1 l_2 \dot{\theta}_2 \cos(\theta_1-\theta_2)
	\]
	\[
	\frac{d}{dt} \frac{\partial L}{\partial \dot{\theta}_1} = (m_1+m_2)l_1^2 \ddot{\theta}_1 + m_2 l_1 l_2 \ddot{\theta}_2 \cos(\theta_1-\theta_2) - m_2 l_1 l_2 \dot{\theta}_2 (\dot{\theta}_1-\dot{\theta}_2) \sin(\theta_1-\theta_2)
	\]
	\[
	\frac{\partial L}{\partial \theta_1} = -m_2 l_1 l_2 \dot{\theta}_1 \dot{\theta}_2 \sin(\theta_1-\theta_2) - (m_1+m_2) g l_1 \sin\theta_1
	\]
	
	方程 $\frac{d}{dt}\frac{\partial L}{\partial \dot{\theta}_1} - \frac{\partial L}{\partial \theta_1} = 0$ 给出：
	\[
	\boxed{(m_1+m_2) l_1 \ddot{\theta}_1 + m_2 l_2 \ddot{\theta}_2 \cos(\theta_1-\theta_2) + m_2 l_2 \dot{\theta}_2^2 \sin(\theta_1-\theta_2) + (m_1+m_2) g \sin\theta_1 = 0}
	\]
	
	\textbf{Lagrange 方程（对 $\theta_2$）} 类似可得：
	\[
	\boxed{m_2 l_2 \ddot{\theta}_2 + m_2 l_1 \ddot{\theta}_1 \cos(\theta_1-\theta_2) - m_2 l_1 \dot{\theta}_1^2 \sin(\theta_1-\theta_2) + m_2 g \sin\theta_2 = 0}
	\]
	
	两式构成非线性耦合常微分方程组。小角度线性化（$\sin\theta \approx \theta,\; \cos\theta \approx 1$，略去速度平方项）后：
	\[
	\begin{bmatrix}
	(m_1+m_2) l_1 & m_2 l_2 \\
	m_2 l_1 & m_2 l_2
	\end{bmatrix}
	\begin{Bmatrix} \ddot{\theta}_1 \\ \ddot{\theta}_2 \end{Bmatrix} + 
	\begin{bmatrix}
	(m_1+m_2) g & 0 \\
	0 & m_2 g
	\end{bmatrix}
	\begin{Bmatrix} \theta_1 \\ \theta_2 \end{Bmatrix} = \mathbf{0}
	\]
	化为标准特征值问题 $\det(\mathbf{K} - \omega^2 \mathbf{M}) = 0$，可解出两阶固有频率 $\omega_1, \omega_2$ 及对应模态。
\end{example}

\section{受非保守力系统的 Lagrange 方程}

实际工程系统中普遍存在阻尼与外激励等非保守力。处理方式如下：

\subsection{广义力的计算}

对于非保守力（给定力 $\mathbf{F}_k^{\text{(nc)}}$ 作用于点 $\mathbf{r}_k$），对应的广义力由定义 (3) 直接计算：
\[
Q_j^{\text{(nc)}} = \sum_{k} \mathbf{F}_k^{\text{(nc)}} \cdot \frac{\partial \mathbf{r}_k}{\partial q_j}
\]
或者通过虚功关系确定：$\delta W^{\text{(nc)}} = \sum_j Q_j^{\text{(nc)}} \delta q_j$。

\subsection{Rayleigh 耗散函数}

对于线性黏性阻尼，引入 \textbf{Rayleigh 耗散函数}：
\begin{myformula}
	R = \frac{1}{2} \sum_{i,j} c_{ij} \dot{q}_i \dot{q}_j
\end{myformula}
其中 $c_{ij}$ 为黏性阻尼系数，$\{c_{ij}\} = \mathbf{C}$ 为半正定对称阻尼矩阵。阻尼力对应的广义力为：
\[
Q_j^{\text{(damp)}} = -\frac{\partial R}{\partial \dot{q}_j} = -\sum_{i} c_{ji} \dot{q}_i
\]

\subsection{推广的 Lagrange 方程}

将阻尼力与外激励纳入 Lagrange 框架：
\begin{myformula}
	\frac{d}{dt} \left( \frac{\partial L}{\partial \dot{q}_j} \right) - \frac{\partial L}{\partial q_j} + \boxed{\frac{\partial R}{\partial \dot{q}_j}} = f_j(t), \quad j = 1, \ldots, n
\end{myformula}
其中 $f_j(t)$ 为外激励对应的广义力（控制力、风载、地震激励等）。

\begin{example}{受迫振动弹簧-阻尼-质量系统}
	质块 $m$、弹簧 $k$、阻尼 $c$，受外激励 $F(t)$。广义坐标 $x$。
	
	动能 $T = \frac{1}{2}m\dot{x}^2$，势能 $V = \frac{1}{2}kx^2$，$L = \frac{1}{2}m\dot{x}^2 - \frac{1}{2}kx^2$。\\
	Rayleigh 函数 $R = \frac{1}{2}c\dot{x}^2$，外激励广义力 $f(t) = F(t)$。
	
	推广 Lagrange 方程：
	\[
	\frac{d}{dt} (m\dot{x}) - (-kx) + c\dot{x} = F(t) \;\Longrightarrow\; m\ddot{x} + c\dot{x} + kx = F(t)
	\]
	这正是经典的受迫振动方程。
\end{example}

\section{Hamilton 原理——积分形式的变分原理}

Hamilton 原理是分析力学最高层次的表述，它将整个运动规律浓缩为一个变分驻值条件。

\subsection{变分法的基本概念}

考虑泛函（函数的函数）：
\[
S[q] = \int_{t_1}^{t_2} L(q(t), \dot{q}(t), t) \, dt
\]
$S[q]$ 为\textbf{作用量泛函}（Action functional），$L$ 为 Lagrange 函数。

变分 $\delta q(t)$ 是连接两条相邻路径的无限小差异，满足端点固定条件：
\[
\delta q(t_1) = \delta q(t_2) = 0
\]

\begin{myTheorem}{Hamilton 原理（最小作用量原理）}
	在给定初末时刻 $t_1, t_2$ 及其对应的位形之间，真实运动 $q(t)$ 使作用量泛函 $S[q]$ 取驻值：
	\[
	\delta S = \delta \int_{t_1}^{t_2} L(q, \dot{q}, t) \, dt = 0
	\]
\end{myTheorem}

\begin{proof}[Hamilton 原理导出 Lagrange 方程]
	考虑含 $n$ 个自由度的系统，$q = (q_1, \ldots, q_n)$。作用量变分：
	\[
	\delta S = \int_{t_1}^{t_2} \sum_{j=1}^{n} \left( \frac{\partial L}{\partial q_j} \delta q_j + \frac{\partial L}{\partial \dot{q}_j} \delta \dot{q}_j \right) dt
	\]
	
	注意到 $\delta \dot{q}_j = \frac{d}{dt}(\delta q_j)$（变分与求导可交换次序），对第二项分部积分：
	\[
	\int_{t_1}^{t_2} \frac{\partial L}{\partial \dot{q}_j} \delta \dot{q}_j \, dt = \left[ \frac{\partial L}{\partial \dot{q}_j} \delta q_j \right]_{t_1}^{t_2} - \int_{t_1}^{t_2} \frac{d}{dt} \left( \frac{\partial L}{\partial \dot{q}_j} \right) \delta q_j \, dt
	\]
	
	由端点条件 $\delta q_j(t_1) = \delta q_j(t_2) = 0$，边界项为零。代入得：
	\[
	\delta S = \int_{t_1}^{t_2} \sum_{j=1}^{n} \left[ \frac{\partial L}{\partial q_j} - \frac{d}{dt} \left( \frac{\partial L}{\partial \dot{q}_j} \right) \right] \delta q_j \, dt = 0
	\]
	
	由 $\delta q_j$ 的任意性（对任意容许变分 $\delta S$ 均为零），被积函数须恒等于零，得\textbf{Euler-Lagrange 方程}：
	\[
	\frac{d}{dt} \left( \frac{\partial L}{\partial \dot{q}_j} \right) - \frac{\partial L}{\partial q_j} = 0, \quad j = 1, \ldots, n
	\]
	这正是无外力的第二类 Lagrange 方程。
\end{proof}

\begin{remark}
	Hamilton 原理揭示了力学规律背后深刻的数学结构——\textbf{变分原理}。Lagrange 方程等价于作用量的 Euler-Lagrange 方程，正如 Fermat 原理（光程最短）等价于几何光学中的 Snell 定律。这种"极值原理→微分方程"的范式广泛存在于物理学的各个分支（电磁学、量子力学、广义相对论等）。
\end{remark}

\subsection{从 Hamilton 原理到守恒律——Noether 定理简介}

Hamilton 原理的一个重要推论是 \textbf{Noether 定理}：作用量的每一个连续对称性对应一个守恒量。具体地：
\begin{itemize}
	\item 时间平移不变性 $\longleftrightarrow$ 能量守恒（$H = T + V$ 为常数）；
	\item 空间平移不变性 $\longleftrightarrow$ 动量守恒；
	\item 空间转动不变性 $\longleftrightarrow$ 角动量守恒。
\end{itemize}

这一结论揭示了物理守恒律与空间-时间对称性之间的深刻联系，是理论物理的基石之一。

\section{Hamilton 力学——Legendre 变换与相空间}

Lagrange 力学在 $n$ 维位形空间中以 $n$ 个二阶微分方程描述运动。\textbf{Hamilton 力学}通过 Legendre 变换将其转化为 $2n$ 个一阶方程，在相空间中描述系统的演化。

\subsection{Legendre 变换与 Hamilton 函数}

定义\textbf{广义动量（共轭动量）}：
\begin{myformula}
	p_j = \frac{\partial L}{\partial \dot{q}_j}, \quad j = 1, \ldots, n
\end{myformula}

Legendre 变换将变量 $(\dot{q}_1, \ldots, \dot{q}_n)$ 替换为 $(p_1, \ldots, p_n)$，定义 \textbf{Hamilton 函数}：
\[
H(q, p, t) = \sum_{j=1}^{n} p_j \dot{q}_j - L(q, \dot{q}, t)
\]
右端的 $\dot{q}_j$ 须由 $p_j = \partial L / \partial \dot{q}_j$ 反解为 $p$ 和 $q$ 的函数。

\begin{example}{Legendre 变换实例——简谐振子}
	$L = \frac{1}{2}m\dot{x}^2 - \frac{1}{2}kx^2$。广义动量 $p = \partial L / \partial \dot{x} = m\dot{x}$，反解得 $\dot{x} = p / m$。
	
	Hamilton 函数：
	\[
	H = p\dot{x} - L = p \cdot \frac{p}{m} - \left( \frac{1}{2}m\frac{p^2}{m^2} - \frac{1}{2}kx^2 \right) = \frac{p^2}{2m} + \frac{1}{2}kx^2 = T + V
	\]
	即系统总机械能。
\end{example}

\subsection{Hamilton 正则方程}

对 $H(q, p, t)$ 求全微分：
\[
dH = \sum_{j} \left( \frac{\partial H}{\partial q_j} dq_j + \frac{\partial H}{\partial p_j} dp_j \right) + \frac{\partial H}{\partial t} dt
\]

另一方面，由 $H = \sum p_j \dot{q}_j - L$：
\[
dH = \sum_j \left( \dot{q}_j dp_j + p_j d\dot{q}_j - \frac{\partial L}{\partial q_j} dq_j - \frac{\partial L}{\partial \dot{q}_j} d\dot{q}_j \right) - \frac{\partial L}{\partial t} dt
\]

由 $p_j = \partial L / \partial \dot{q}_j$，$p_j d\dot{q}_j$ 与 $-\frac{\partial L}{\partial \dot{q}_j} d\dot{q}_j$ 抵消。又由 Lagrange 方程 $\dot{p}_j = \frac{d}{dt}(\partial L/\partial \dot{q}_j) = \partial L/\partial q_j$，比较系数得 \textbf{Hamilton 正则方程}：
\[
\boxed{\dot{q}_j = \frac{\partial H}{\partial p_j}, \qquad \dot{p}_j = -\frac{\partial H}{\partial q_j}}, \quad j = 1, \ldots, n
\]

这组 $2n$ 个一阶方程具有优美的对称结构，且 $H$ 沿真实运动的时间导数为：
\[
\frac{dH}{dt} = \sum_j \left( \frac{\partial H}{\partial q_j} \dot{q}_j + \frac{\partial H}{\partial p_j} \dot{p}_j \right) + \frac{\partial H}{\partial t} = \frac{\partial H}{\partial t}
\]
若 $H$ 不显含时间（$\partial H / \partial t = 0$），则 $H$ 守恒。

\subsection{辛几何结构与数值积分的关联}

Hamilton 方程的流保持相空间中的\textbf{辛 2-形式} $\omega = \sum_j dp_j \wedge dq_j$ 不变，即相体积（Liouville 定理）和 Poincar\'{e} 积分不变量守恒。这一几何性质在现代数值分析中具有重要应用。

\textbf{辛几何积分器}（Symplectic integrator）是专门设计用于保持 Hamilton 系统辛结构的数值方法，相较于一般的 Runge-Kutta 方法，它们在长时间能量行为和相空间结构保持方面具有显著优势。

经典的 \textbf{St\"{o}rmer-Verlet} 方法（速度 Verlet 格式）：
\begin{myformula}
	\begin{cases}
		\mathbf{p}_{n+1/2} = \mathbf{p}_n - \dfrac{h}{2} \nabla_q V(\mathbf{q}_n) \\[8pt]
		\mathbf{q}_{n+1} = \mathbf{q}_n + h M^{-1} \mathbf{p}_{n+1/2} \\[8pt]
		\mathbf{p}_{n+1} = \mathbf{p}_{n+1/2} - \dfrac{h}{2} \nabla_q V(\mathbf{q}_{n+1})
	\end{cases}
\end{myformula}
其构造采用了"力-位置-力"的三段式分裂策略，本质上是对 Hamilton 向量场的对称 Trotter 分解，保证了方法的时间可逆性与二阶精度。

\subsection{能量法在工程建模中的统一视角}

回顾本章与前述各章：《数值解法》提供了求解常微分方程的通用工具，《单/多自由度系统》建立了二阶振动方程的标准形式，《状态空间方法》引入了系统的一阶等价表示。而\textbf{Lagrange-Hamilton 体系}提供了从物理原理自动生成这些方程的统一框架：

\begin{enumerate}
	\item \textbf{Lagrange 层面}：输入 $T$、$V$ 和耗散函数 $R$，自动输出 $\mathbf{M}\ddot{\mathbf{q}} + \mathbf{C}\dot{\mathbf{q}} + \mathbf{K}\mathbf{q} = \mathbf{f}$；
	\item \textbf{Hamilton 层面}：自动获得一阶状态空间形式 $\dot{\mathbf{x}} = \mathbf{A}\mathbf{x} + \mathbf{B}\mathbf{u}$；
	\item \textbf{数值积分}：可选择通用 RK 方法或专门的辛积分器进行时间推进。
\end{enumerate}

\begin{remark}
	在工程动力学仿真软件（如 ADAMS、Simpack、Simscape Multibody 等）中，多体系统的运动方程正是通过 Lagrange 方法（或以之为基础的递推算法）自动组集质量矩阵和广义力向量，再配合数值积分器求解，实现从 CAD 几何模型到动力学响应的一体化自动计算。
\end{remark}
